%% kapa: exam=2 round=29 year=2018 period=3 subject=compensation_law subject_ko=감정평가 및 보상법규 seq=1 number=1 points=40
%% source: Q-Net 감정평가사 기출문제 (cst00309) artlSeq=5203398 file=raw/1008/5203398/3교시 감정평가 및 보상법규.pdf
%% conversion: kordoc 4.24.51 -> tools/gen2.py (auto); review=proofread-gpt-terra; flags=none
\begin{longquestion}{【문제 1】}
A도 도지사 甲은 도내의 심각한 주차난을 해결하기 위하여 A도내 B시 일대 40,000 m\textsuperscript{2}(이하 ‘이 사건 공익사업구역’이라 함)를 공영주차장으로 사용하고자 사업계획을 수립하고 「공익사업을 위한 토지 등의 취득 및 보상에 관한 법률」(이하 ‘토지보상법’이라 함)에 따른 절차를 거쳐, 국토교통부장관의 사업인정을 받고 이를 고시하였다. 이후 甲은 이 사건 공익사업구역 내 주택 세입자 乙 등이 이 사건 공익사업이 시행되는 동안 임시로 거주할 수 있도록 B시에 임대아파트를 건립하여 세입자에게 제공하는 등 이주대책을 수립ㆍ시행하였다. 한편, 乙은 「공익사업을 위한 토지 등의 취득 및 보상에 관한 법률 시행규칙」(이하 ‘토지보상법 시행규칙’이라 함) 제54조 제2항에 해당하는 세입자이다. 다음 물음에 답하시오. (40점) \par
1) 乙은 토지보상법 시행규칙에 따른 주거이전비를 받을 수 있는 권리를 포기한다는 취지의 ‘임대아파트 입주에 따른 주거이전비 포기각서’를 甲에게 제출하고 위 임대아파트에 입주하였지만, 이후 관련 법령이 임대아파트와 같은 임시수용시설 등을 제공받는 자를 주거이전비 지급대상에서 배제하지 않고 있는 점을 알게 되었다. 이에 乙은 위 포기각서를 무시하고 토지보상법 시행규칙 상의 주거이전비를 청구하였다. 乙의 주거이전비 청구의 인용여부에 관하여 논하시오. (30점) \par
2) 한편, 丙은 이 사건 공익사업구역 밖에서 음식점을 경영하고 있었는데, 이 사건 공익사업으로 인하여 자신의 음식점의 주출입로가 단절되어 일정 기간 휴업을 할 수 밖에 없게 되었다. 이때, 丙은 토지보상법령상 보상을 받을 수 있는가? (10점) \par
\begin{framedtext}
참조 조문 \par
<공익사업을 위한 토지 등의 취득 및 보상에 관한 법률 시행규칙> 
제54조(주거이전비의 보상) ① 공익사업시행지구에 편입되는 주거용 건축물의 소유자에 대하여는 해당 건축물에 대한 보상을 하는 때에 가구원수에 따라 2개월분의 주거이전비를 보상하여야 한다. 다만, 건축물의 소유자가 해당 건축물 또는 공익사업시행지구 내 타인의 건축물에 실제 거주하고 있지 아니하거나 해당 건축물이 무허가건축물등인 경우에는 그러하지 아니하다. 
② 공익사업의 시행으로 인하여 이주하게 되는 주거용 건축물의 세입자(법 
제78조제1항에 따른 이주대책대상자인 세입자는 제외한다)로서 사업인정고시일등 당시 또는 공익사업을 위한 관계법령에 의한 고시 등이 있은 당시 해당 공익사업시행지구안에서 3개월 이상 거주한 자에 대하여는 가구원수에 따라 4개월분의 주거이전비를 보상하여야 한다. 다만, 무허가건축물등에 입주한 세입자로서 사업인정고시일등 당시 또는 공익사업을 위한 관계법령에 의한 고시 등이 있은 당시 그 공익사업지구 안에서 1년 이상 거주한 세입자에 대하여는 본문에 따라 주거이전비를 보상하여야 한다.
\end{framedtext}

\end{longquestion}
